\documentclass[11pt]{article}
\usepackage[a4paper,margin=1in]{geometry}
\usepackage{fontspec}
\usepackage[boldfont=false]{xeCJK}
\setCJKmainfont{FandolSong-Regular.otf}
\usepackage{amsmath}
\usepackage{amssymb}
\usepackage{array}
\usepackage{booktabs}
\usepackage{graphicx}
\usepackage{adjustbox}
\usepackage{enumitem}
\setlist[itemize]{topsep=2pt,partopsep=0pt,parsep=0pt,itemsep=2pt,leftmargin=1.4em}
\setlist[enumerate]{topsep=2pt,partopsep=0pt,parsep=0pt,itemsep=2pt,leftmargin=1.6em}
\usepackage{parskip}
\usepackage{fvextra}
\fvset{breaklines=true,breakanywhere=true,fontsize=\small}
\setlength{\parindent}{0pt}

\begin{document}

\begin{center}
  {\LARGE\bfseries Variation and selection}\\[0.35em]
  {\large Igcse Biology}
\end{center}
\vspace{0.6em}

\section*{Variation}

\textbf{Variation} 变异 means the differences between individuals of the same \textbf{species} 物种. There are two types:

\begin{itemize}
\item \textbf{continuous variation} 连续变异 — a smooth range of \textbf{phenotypes} 表现型 between two extremes (for example body length and body mass).

\item \textbf{discontinuous variation} 不连续变异 — a few separate phenotypes with nothing in between (for example \textbf{ABO blood groups} 血型, or seed shape in peas).

\end{itemize}

Discontinuous variation is usually caused by \textbf{genes} 基因 only. Continuous variation is caused by \textbf{both} genes and the \textbf{environment} 环境 — for example, your height depends on your genes and also on your diet.

\subsection*{Mutation}

A \textbf{mutation} 突变 is a genetic change. Mutations are the way \textbf{new} \textbf{alleles} 等位基因 are made, so they are the original source of all variation. \textbf{(Supplement)} A \textbf{gene mutation} 基因突变 is a random change in the \textbf{base sequence} 碱基序列 of DNA.

\textbf{Ionising radiation} 电离辐射 (such as X-rays) and some chemicals raise the \textbf{rate} 速率 of mutation. \textbf{(Supplement)} Other sources of genetic variation are \textbf{meiosis} 减数分裂, random mating and random \textbf{fertilisation} 受精.

\section*{Adaptive features}

An \textbf{adaptive feature} 适应特征 is an \textbf{inherited} 遗传 feature that helps an organism to \textbf{survive} 生存 and \textbf{reproduce} 生殖 in its \textbf{environment}.

\textbf{(Supplement)}

\begin{itemize}
\item \textbf{hydrophytes} 水生植物 (water plants) have features such as air spaces to help them float, and stomata on the upper leaf surface.

\item \textbf{xerophytes} 旱生植物 (desert plants) have features such as a thick waxy cuticle, few small stomata, and the ability to store water — all to reduce water loss.

\end{itemize}

\section*{Natural selection}

\textbf{Natural selection} 自然选择 explains how a species becomes better suited to its environment:

\begin{enumerate}
\item there is genetic variation in a population (caused by mutation).

\item organisms produce \textbf{many offspring} 后代 — more than can survive.

\item there is a struggle to survive, with \textbf{competition} 竞争 for \textbf{resources} 资源 such as food and space.

\item the individuals that are better \textbf{adapted} 适应 are more likely to survive and reproduce.

\item they pass on their alleles, so the helpful alleles become more common in the next generation.

\end{enumerate}

\textbf{(Supplement)} Over many generations this makes the population more suited to its environment; this process is called adaptation. A clear example is \textbf{antibiotic} 抗生素-resistant \textbf{bacteria} 细菌: a few bacteria carry a \textbf{resistant} 耐药 allele and survive the antibiotic, then multiply, until the whole population is resistant.

\section*{Selective breeding}

In \textbf{selective breeding} 选择育种 (also called \textbf{artificial selection} 人工选择), humans choose which organisms breed:

\begin{enumerate}
\item choose the individuals that have the features you want.

\item cross them to produce the next generation.

\item choose the offspring that show those features, and repeat over many generations.

\end{enumerate}

This is used to improve \textbf{crops} 农作物 (for example a higher yield) and farm animals (for example cows that give more milk).

\textbf{(Supplement)} The key difference: in natural selection the \textbf{environment} decides which individuals survive; in artificial selection \textbf{humans} decide.

\section*{Exam tips}

\begin{itemize}
\item Continuous variation = a range (genes + environment); discontinuous = separate groups (genes only).

\item Mutation makes new alleles — the source of all variation. Radiation and some chemicals raise the mutation rate.

\item Natural selection: variation → too many offspring → struggle and competition → the best-adapted survive and reproduce → they pass on their alleles.

\item Antibiotic resistance in bacteria is natural selection in action.

\item Selective breeding (artificial selection): humans pick the parents; in natural selection, the environment "picks".

\end{itemize}


\end{document}
